\notepage{56}{The Coulomb virial relation and the mean inverse radius}
For the Coulomb potential, take $k=-1$ and $c=-e^2$. The virial theorem
becomes
\[
 \left\langle\frac{\mathbf P^2}{2m}\right\rangle
 =-\frac12\left\langle-\frac{e^2}{R}\right\rangle.
\]
For a normalized hydrogenic energy eigenstate $\ket{n,\ell,m}$,
\begin{align*}
 E_n
 &=\left\langle\frac{\mathbf P^2}{2m}-\frac{e^2}{R}\right\rangle_{n\ell m}\\
 &=\langle T\rangle_{n\ell m}+\langle V\rangle_{n\ell m}\\
 &=\frac12\langle V\rangle_{n\ell m}
 =-\frac{e^2}{2}\left\langle\frac1R\right\rangle_{n\ell m}.
\end{align*}
On the other hand,
\[
 E_n=-\frac{me^4}{2\hbar^2n^2}
 =-\frac{e^2}{2a_0n^2}.
\]
Comparing the two expressions gives
\[
 \boxed{\left\langle\frac1R\right\rangle_{n\ell m}
 =\frac1{n^2a_0}.}
\]
This expectation is independent of $\ell$ and $m$ in the
nonrelativistic Coulomb problem considered here.
\editorial{The potential contribution to $H$ is $+V=-e^2/R$.
The intermediate handwritten expression that subtracts
$\langle-e^2/R\rangle$ is a sign slip; the final result in the notes is
consistent with the corrected addition above.}
